\documentclass[12pt,a4paper]{article}

\usepackage[T1]{fontenc}
\usepackage[utf8]{inputenc}
\usepackage[brazil]{babel}
\usepackage{lmodern}
\usepackage{geometry}
\usepackage{booktabs}
\usepackage{array}
\usepackage{hyperref}
\usepackage{enumitem}
\usepackage{listings}
\usepackage{xcolor}

\geometry{margin=2.5cm}
\hypersetup{
    colorlinks=true,
    linkcolor=black,
    urlcolor=blue
}

\lstset{
    basicstyle=\ttfamily\small,
    breaklines=true,
    frame=single,
    columns=fullflexible
}

\title{Implementação de um TAD Grafo para Escolha de Galpão Logístico}
\author{João CRM}
\date{Junho de 2026}

\begin{document}

\maketitle

\begin{abstract}
Este trabalho apresenta uma implementação em C++ de um TAD de grafo para apoiar a escolha de um bairro para instalação de um galpão logístico. A solução lê os dados de vértices e arestas a partir de arquivos texto, calcula distâncias mínimas com os algoritmos de Dijkstra e Bellman-Ford, estima o custo operacional mensal e anual de cada candidato e escolhe o melhor bairro com base na expressão custoGalpão = custoInstalação – custoOperaçãoAnual. Também foi incluída uma rotina de benchmark que grava em CSV as médias de tempo dos dois algoritmos após várias repetições.
\end{abstract}

\section{Introdução}

A problemática do trabalho consiste em encontrar o melhor bairro para a instalação de um galpão logístico. Para isso, cada bairro possui um custo fixo de instalação, uma demanda estimada para cada mês do ano e conexões com outros bairros vizinhos, representadas por arestas com custo de transporte.

Como o custo de entrega é proporcional à distância entre o bairro onde o galpão será instalado e os demais bairros atendidos, a solução depende do cálculo das menores distâncias dentro da cidade. Dessa forma, foram utilizados dois algoritmos clássicos de caminhos mínimos:

\begin{itemize}[noitemsep]
\item Dijkstra;
\item Bellman-Ford.
\end{itemize}

Embora os dois algoritmos tenham como objetivo encontrar as mesmas menores distâncias, eles apresentam custos computacionais diferentes. Por esse motivo, além de resolver o problema proposto, também é possível comparar o desempenho de cada abordagem.


\section{Estrutura da Solução}

Os arquivos principais da implementação são:

\begin{itemize}[noitemsep]
    \item \texttt{src/Grafo.h}: declara as estruturas \texttt{Bairro}, \texttt{Aresta}, \texttt{ResultadoAlgoritmo}, \texttt{ResultadoAvaliacaoBairro}, \texttt{ResultadoEscolhaGalpao}, \texttt{ResultadoBenchmarkBairro} e a classe \texttt{Grafo};
    \item \texttt{src/Grafo.cpp}: implementa o TAD e os algoritmos;
    \item \texttt{src/main.cpp}: contém o programa principal e o modo de benchmark;
    \item \texttt{dados/vertices.txt}: arquivo de teste com os bairros;
    \item \texttt{dados/arestas.txt}: arquivo de teste com as conexões;
    \item \texttt{dados/vertices\_maior.txt}: segunda instância de teste com mais vértices;
    \item \texttt{dados/arestas\_maior.txt}: segunda instância de teste para o benchmark ampliado;
    \item \texttt{resultados/time\_compare.csv}: saída do benchmark;
    \item \texttt{resultados/time\_compare\_large.csv}: saída do benchmark ampliado;
    \item \texttt{Makefile}: compilação e execução.
\end{itemize}

\subsection{Modelagem adotada}

O grafo foi modelado como não direcionado, assumindo que o custo de deslocamento entre dois bairros é o mesmo nos dois sentidos. Com isso, cada conexão é armazenada de forma simétrica na estrutura de adjacência, o que simplifica a representação.

Na modelagem orientada a objetos, a classe \texttt{Grafo} concentra a estrutura de dados e as operações principais da solução. Ela armazena os bairros e as arestas, executa os algoritmos de caminhos mínimos, calcula os custos operacionais e reúne os resultados da análise. Já o programa principal fica responsável apenas por carregar os arquivos, acionar a execução e exibir as saídas. Essa separação procura manter responsabilidades mais bem definidas entre os módulos, em linha com princípios de coesão e responsabilidade única discutidos por Martin \cite{martin2009,martin2018}.


Cada vértice armazena:

\begin{itemize}[noitemsep]
    \item identificador do bairro;
    \item custo de instalação do galpão;
    \item vetor com 12 demandas mensais.
\end{itemize}

Cada aresta armazena:

\begin{itemize}[noitemsep]
    \item vértice de origem;
    \item vértice de destino;
    \item custo unitário de transporte.
\end{itemize}

Para simplificar a implementação e manter boa eficiência no Dijkstra, foram usadas simultaneamente uma lista de arestas e uma lista de adjacência.

\section{Implementação do TAD}

A implementação foi desenvolvida em C++, seguindo a abordagem da linguagem e de suas bibliotecas padrão apresentada por Stroustrup \cite{stroustrup2013}.

Além do encapsulamento dos dados internos do grafo, a modelagem orientada a objetos foi reforçada ao concentrar na própria classe \texttt{Grafo} as responsabilidades de avaliar cada bairro candidato, escolher o melhor local para instalação e produzir os resultados do benchmark. Com isso, o arquivo \texttt{main.cpp} ficou restrito à apresentação dos resultados e ao tratamento dos argumentos de linha de comando.

Essa organização evita espalhar a lógica principal da solução por funções externas e deixa a orquestração de entrada e saída separada do processamento central. A escolha segue a ideia de manter cada parte do programa com um papel mais claro e com menor acoplamento estrutural, como defendido em \textit{Clean Code} e \textit{Clean Architecture} \cite{martin2009,martin2018}.

\subsection{Leitura dos arquivos}

O método \texttt{carregarVertices} lê a quantidade de vértices e, em seguida, uma linha por bairro contendo identificador, custo de instalação e as 12 demandas mensais.

O método \texttt{carregarArestas} lê a quantidade de arestas e, em seguida, uma linha por aresta contendo os dois identificadores de bairros conectados e o custo da aresta.

Se algum arquivo estiver inválido, a implementação interrompe a execução com uma exceção e mensagem descritiva.

\subsection{Cálculo das distâncias}

O TAD oferece os métodos:

\begin{itemize}[noitemsep]
    \item \texttt{dijkstra(int indiceOrigem)};
    \item \texttt{bellmanFord(int indiceOrigem)}.
\end{itemize}

Cada método retorna:

\begin{itemize}[noitemsep]
    \item o vetor de distâncias mínimas;
    \item o tempo de execução medido em microssegundos.
\end{itemize}

Os vetores retornados pelos dois algoritmos são comparados. Caso os resultados sejam diferentes, a execução é encerrada, pois isso indicaria erro de implementação.

Para tornar a modelagem mais alinha aos princípios da orientação a objetos, a análise de negócio também foi incorporada ao TAD por meio dos métodos:

\begin{itemize}[noitemsep]
    \item \texttt{avaliarBairroCandidato(int indiceDeposito)};
    \item \texttt{escolherMelhorGalpao()};
    \item \texttt{executarBenchmark(int repeticoes)}.
\end{itemize}

Esses métodos retornam objetos de resultado que agrupam as informações calculadas, evitando que a lógica de avaliação e decisão fique espalhada fora da classe principal do domínio.

\subsection{Cálculo do custo operacional}

Para cada bairro candidato a depósito:

\begin{enumerate}[noitemsep]
    \item calcula-se a distância mínima até todos os demais bairros;
    \item multiplica-se cada distância pela demanda do bairro em um dado mês;
    \item soma-se o custo de todos os bairros para obter o custo operacional mensal;
    \item soma-se o custo dos 12 meses para obter o custo operacional anual.
\end{enumerate}


Ao final, a decisão é representada pela expressão:

\begin{lstlisting}
custoGalpao = custoInstalacao - custoOperacaoAnual
\end{lstlisting}

Essa expressão pode ser interpretada não como um custo total no sentido estrito, mas como um saldo entre o investimento fixo inicial e o impacto anual da operação. Nessa leitura, valores negativos indicam que o custo operacional anual superou o custo fixo de instalação, o que ajuda a explicar os resultados obtidos. Assim, o critério adotado compara diretamente essas duas parcelas para definir o bairro escolhido.


\section{Análise de Complexidade}

Esta seção considera apenas os métodos do TAD.

\subsection{\texttt{carregarVertices}}

\begin{itemize}[noitemsep]
    \item Tempo: $O(V)$;
    \item Espaço adicional: $O(V)$.
\end{itemize}

\subsection{\texttt{carregarArestas}}

\begin{itemize}[noitemsep]
    \item Tempo: $O(E \cdot V)$;
    \item Espaço adicional: $O(E + V)$.
\end{itemize}

Esse custo ocorre porque a implementação simples faz a busca do índice de cada identificador por varredura linear no vetor de bairros.

\subsection{\texttt{indicePorId}}

\begin{itemize}[noitemsep]
    \item Tempo: $O(V)$;
    \item Espaço adicional: $O(1)$.
\end{itemize}

\subsection{\texttt{dijkstra}}

\begin{itemize}[noitemsep]
    \item Tempo: $O((V + E)\log V)$;
    \item Espaço adicional: $O(V)$.
\end{itemize}

\subsection{\texttt{bellmanFord}}

\begin{itemize}[noitemsep]
    \item Tempo: $O(V \cdot E)$;
    \item Espaço adicional: $O(V)$.
\end{itemize}

\subsection{\texttt{calcularCustoOperacaoMensal}}

\begin{itemize}[noitemsep]
    \item Tempo: $O(V)$;
    \item Espaço adicional: $O(1)$.
\end{itemize}

\subsection{\texttt{calcularCustoOperacaoAnual}}

\begin{itemize}[noitemsep]
    \item Tempo: $O(12 \cdot V)$, equivalente a $O(V)$;
    \item Espaço adicional: $O(1)$.
\end{itemize}

\section{Dados de Teste}

Foram usadas duas instâncias de teste. A primeira possui 5 bairros, a segunda foi adicionada apenas para fortalecer a comparação experimental entre os algoritmos.

\subsection{Arquivo de vértices}

\begin{lstlisting}
5
1 1200 10 12 11 13 14 16 15 14 13 12 11 10
2 950 8 9 10 10 11 12 13 14 13 12 10 9
3 1100 15 14 16 17 18 18 17 16 15 14 14 15
4 980 6 7 8 9 10 11 10 9 8 7 6 6
5 1050 9 10 10 11 12 13 14 14 13 12 11 10
\end{lstlisting}

\subsection{Arquivo de arestas}

\begin{lstlisting}
7
1 2 4.5
1 3 2.0
2 3 1.5
2 4 3.0
3 4 2.5
3 5 4.0
4 5 1.0
\end{lstlisting}

O grafo de teste é conexo.

\subsection{Segunda instância para benchmark}

Para ampliar a análise experimental, também foi criada uma instância com 20 bairros e 42 arestas, armazenada nos arquivos \texttt{dados/vertices\_maior.txt} e \texttt{dados/arestas\_maior.txt}. Essa instância continua sendo conexa e usa apenas pesos positivos, preservando as premissas do problema.

\section{Resultados}

\subsection{Custos anuais calculados}

\begin{table}[h]
\centering
\begin{tabular}{>{\raggedright\arraybackslash}p{2cm}rrr}
\toprule
Bairro & Custo de instalação & Custo operacional anual & \texttt{custoGalpao} \\
\midrule
1 & 1200,00 & 2037,50 & -837,50 \\
2 & 950,00 & 1659,00 & -709,00 \\
3 & 1100,00 & 1227,50 & -127,50 \\
4 & 980,00 & 1684,00 & -704,00 \\
5 & 1050,00 & 2113,00 & -1063,00 \\
\bottomrule
\end{tabular}
\caption{Resultados obtidos com os dados de teste.}
\end{table}

Observa-se que todos os valores de \texttt{custoGalpao} ficaram negativos, pois em todas as instâncias testadas o custo operacional anual superou o custo fixo de instalação. Dentro da interpretação adotada neste trabalho, isso pode ser visto como um saldo líquido negativo entre o investimento inicial e a operação anual.

Pela regra, o melhor bairro foi o bairro 5, pois apresentou o menor valor de \texttt{custoGalpao}.

\subsection{Comparação de tempos}

Antes dos resultados experimentais, a Tabela~\ref{tab:comparacao-teorica} resume as principais diferenças conceituais entre os algoritmos usados.

\begin{table}[h]
\centering
\begin{tabular}{p{3.1cm}p{4.2cm}p{4.2cm}}
\toprule
Critério & Dijkstra & Bellman-Ford \\
\midrule
Premissa sobre pesos & Exige pesos não negativos para garantir corretude. & Aceita pesos negativos, desde que não haja ciclo negativo alcançável. \\
Complexidade nesta implementação & $O((V + E)\log V)$ com fila de prioridade e lista de adjacência. & $O(V \cdot E)$ com relaxamento repetido da lista de arestas. \\
Estrutura principal usada & Lista de adjacência e fila de prioridade. & Lista de arestas percorrida em várias iterações. \\
Vantagem esperada & Melhor escalabilidade em grafos maiores com pesos não negativos. & Implementação simples e robusta para casos mais gerais. \\
\bottomrule
\end{tabular}
\caption{Comparação teórica entre Dijkstra e Bellman-Ford.}
\label{tab:comparacao-teorica}
\end{table}

Além da execução normal, foi criado um modo de benchmark no programa para repetir os algoritmos várias vezes por bairro e gravar as médias em CSV. O alvo:

\begin{lstlisting}
make run-time-compare
\end{lstlisting}

gera o arquivo:

\begin{lstlisting}
resultados/time_compare.csv
\end{lstlisting}

Para a segunda instância, o alvo:

\begin{lstlisting}
make run-time-compare-large
\end{lstlisting}

gera:

\begin{lstlisting}
resultados/time_compare_large.csv
\end{lstlisting}

Esse arquivo segue o formato:

\begin{lstlisting}
bairro_id,repeticoes,media_dijkstra_us,media_bellman_ford_us
\end{lstlisting}

A medição do benchmark é feita externamente às funções do TAD, acumulando o tempo real de cada chamada em ponto flutuante. Isso evita que, em entradas pequenas, a conversão imediata para microssegundos inteiros zere artificialmente a média.

Com $100000$ repetições por bairro, foi obtido o CSV abaixo:

\begin{table}[h]
\centering
\begin{tabular}{crr}
\toprule
Bairro & Média Dijkstra (us) & Média Bellman-Ford (us) \\
\midrule
1 & 0,2331 & 0,1016 \\
2 & 0,2103 & 0,0990 \\
3 & 0,1947 & 0,0880 \\
4 & 0,2017 & 0,1007 \\
5 & 0,1967 & 0,1007 \\
\bottomrule
\end{tabular}
\caption{Médias medidas no benchmark para os dados de teste.}
\end{table}

Na instância mínima, Bellman-Ford foi mais rápido em todos os bairros testados. A média geral ficou em $0{,}0980$ us contra $0{,}2073$ us do Dijkstra, isto é, Dijkstra levou em média 2,12 vezes o tempo do Bellman-Ford. Como a razão variou apenas de 1,95 a 2,29 entre os bairros, o comportamento foi consistente dentro dessa entrada.

Para reduzir o risco de tirar conclusão a partir de um único caso muito pequeno, a segunda instância com 20 bairros também foi submetida ao benchmark com $20000$ repetições por bairro. Em vez de listar as 20 linhas do CSV, a Tabela~\ref{tab:resumo-benchmark} resume as médias gerais das duas instâncias.

\begin{table}[h]
\centering
\begin{tabular}{lrrr}
\toprule
Instância & Média Dijkstra (us) & Média Bellman-Ford (us) & Razão Dijkstra/Bellman \\
\midrule
5 bairros, 7 arestas & 0,2073 & 0,0980 & 2,12 \\
20 bairros, 42 arestas & 0,4902 & 0,3125 & 1,59 \\
\bottomrule
\end{tabular}
\caption{Resumo das médias gerais dos benchmarks executados.}
\label{tab:resumo-benchmark}
\end{table}

Na instância maior, Bellman-Ford continuou mais rápido em todos os vértices medidos. A média geral passou para $0{,}3125$ us, enquanto Dijkstra ficou em $0{,}4902$ us. A razão média caiu para 1,59, com variação entre 1,29 e 1,93. Isso é relevante porque mostra que, embora Bellman-Ford ainda tenha vencido, a vantagem relativa diminuiu quando o grafo cresceu.

Os resultados experimentais, portanto, não repetiram a expectativa puramente assintótica. A explicação mais plausível é que as entradas avaliadas ainda são pequenas para destacar a vantagem estrutural do Dijkstra, enquanto Bellman-Ford se beneficia aqui de uma implementação simples, com percursos lineares sobre um vetor de arestas e pouca sobrecarga de estruturas auxiliares. Já o Dijkstra, apesar da melhor ordem assintótica, paga o custo constante da fila de prioridade mesmo em grafos modestos.

Assim, a comparação entre os algoritmos pode ser resumida em três pontos. Primeiro, ambos produziram os mesmos caminhos mínimos em todas as execuções, confirmando a corretude esperada. Segundo, Bellman-Ford foi mais rápido nas duas instâncias avaliadas. Terceiro, o aumento do tamanho do grafo reduziu a vantagem relativa de Bellman-Ford, o que é compatível com a hipótese de que, em entradas substancialmente maiores, a tendência teórica do Dijkstra pode começar a aparecer de forma mais clara.

\section{Compilação e Execução}

Compilação:

\begin{lstlisting}
make
\end{lstlisting}

Execução principal:

\begin{lstlisting}
make run
\end{lstlisting}

Benchmark com CSV:

\begin{lstlisting}
make run-time-compare
\end{lstlisting}

Benchmark da instância maior:

\begin{lstlisting}
make run-time-compare-large
\end{lstlisting}

\section{Conclusão}

O trabalho implementa um TAD de grafo em C++ capaz de modelar os bairros da cidade, calcular caminhos mínimos com Dijkstra e Bellman-Ford e estimar o custo operacional anual de um galpão logístico. A solução também compara os algoritmos por meio de benchmarks repetidos em duas instâncias, com saída em CSV, o que fortalece a análise experimental do relatório.

Sob a perspectiva de orientação a objetos, a solução também procurou manter encapsulamento, coesão e delegação de responsabilidades. A classe \texttt{Grafo} concentra os dados e os comportamentos centrais do problema, enquanto o programa principal atua apenas como interface de execução. Essa separação contribui para uma modelagem mais consistente e facilita a evolução incremental do código.

Como extensões futuras, seria possível:

\begin{itemize}[noitemsep]
    \item adicionar uma estrutura de mapeamento direto de identificadores para índices;
    \item validar automaticamente se o grafo de entrada é conexo;
    \item incluir mais instâncias de teste para observar melhor o comportamento assintótico dos algoritmos.
\end{itemize}

\begin{thebibliography}{9}

\bibitem{stroustrup2013}
Bjarne Stroustrup.
\textit{The C++ Programming Language}.
4th ed.
Addison-Wesley, 2013.

\bibitem{martin2009}
Robert C. Martin.
\textit{Clean Code: A Handbook of Agile Software Craftsmanship}.
Prentice Hall, 2009.

\bibitem{martin2018}
Robert C. Martin.
\textit{Clean Architecture: A Craftsman's Guide to Software Structure and Design}.
Prentice Hall, 2018.

\end{thebibliography}

\end{document}
